\chapter{Movimiento Lateral y Correlación Multisistema}

En entornos de cómputo académicos e institucionales, los servidores rara vez operan como silos aislados. La integración funcional de servicios suele apoyarse en esquemas de confianza mutua: montajes remotos de directorios mediante el protocolo NFS (\textit{Network File System}), unificación de nombres de usuario y bases de datos de contraseñas compartidas mediante NIS o sincronización periódica, y permisos de acceso simplificados entre equipos de una misma subred física.

Esta interconexión, sumamente ventajosa para la productividad de los usuarios, se transforma en un vector crítico de propagación cuando un sistema es comprometido. Este capítulo expone el análisis pericial conducido sobre el servidor secundario \texttt{victima2} (\texttt{192.0.2.20}), documentando la correlación criptográfica de artefactos idénticos hallados en múltiples cuentas y la reconstrucción del movimiento lateral del atacante.

\section{Arquitectura Departamental y Vectores de Confianza}

Dentro de la topología departamental, \texttt{victima1} y \texttt{victima2} operaban en estrecha sinergia:
\begin{itemize}
    \item Los usuarios del departamento contaban con cuentas homologadas en ambos servidores, manteniendo los mismos nombres de inicio de sesión (\textit{login}) e idénticos identificadores numéricos de usuario y grupo (UID/GID).
    \item Ambas máquinas compartían directorios de datos y volúmenes de almacenamiento temporal montados vía NFS sobre la red local Fast Ethernet.
    \item Los usuarios saltaban con frecuencia de una máquina a otra mediante conexiones interactivas SSH o comandos remotos (\texttt{rsh}/\texttt{ssh}) sin fricción operativa.
\end{itemize}

Al confirmarse la intrusión en \texttt{victima1}, el protocolo pericial exigió auditar de inmediato la imagen forense de \texttt{victima2} para determinar si el compromiso se encontraba confinado o si había derivado en una infección en cascada.

\section{Inspección Forense de victima2: Artefactos Espejo}

El examen del directorio personal del usuario \texttt{duende} en el servidor secundario \texttt{victima2} arrojó un hallazgo determinante: la existencia de paquetes de herramientas idénticos a los descubiertos en la máquina primaria.

En la ruta \texttt{/home/staff/duende/test.tar} se localizó un archivo comprimido de 61,440 bytes. Su desglose interno se detalla en el Listado~\ref{lst:test_tar_victima2}.

\begin{lstlisting}[style=terminal,caption={Contenido del archivo test.tar localizado en victima2:/home/staff/duende.},label={lst:test_tar_victima2}]
drwx------ scott/users       0 2002-02-08 07:09:51 ex/
-rwx------ scott/users   14361 2002-02-08 07:09:35 ex/ep
-rwx------ scott/users   14543 2002-02-08 07:09:39 ex/rh
-rwx------ scott/users   23577 2002-02-08 07:09:45 ex/sd
\end{lstlisting}

\subsection{Cotejo Criptográfico y Vínculo con Sendmail}

El examen funcional y criptográfico de este paquete confirmó que:
\begin{enumerate}
    \item Los binarios \texttt{ex/ep}, \texttt{ex/rh} y \texttt{ex/sd} correspondían a programas de explotación precompilados (\textit{exploits}) dirigidos contra vulnerabilidades en versiones históricas de \texttt{Sendmail} en diferentes arquitecturas de Linux (Red Hat, Slackware y Debian).
    \item Este paquete era **bit a bit idéntico** al localizado en los directorios de las cuentas \texttt{zena} y \texttt{proa} en el servidor \texttt{victima1}.
    \item En el subdirectorio temporal \texttt{temp/} de la cuenta \texttt{duende} en \texttt{victima2} se identificaron los bloques correspondientes a la estructura eliminada \texttt{ex/}, demostrando que el usuario desempaquetó los ejecutables, ejecutó pruebas y posteriormente intentó suprimir el directorio para no dejar rastros.
\end{enumerate}

Asimismo, se recuperó en el espacio no asignado un archivo comprimido suprimido denominado \texttt{o0o.tar}. El cálculo de su suma de verificación criptográfica arrojó un hash MD5 exactamente igual al del archivo homónimo descubierto en \texttt{victima1}, probando la transferencia y replicación manual de las herramientas entre ambos servidores.

\section{Correlación de Accesos Remotos y Direccionamiento IP}

La auditoría de las bitácoras de inicio de sesión (\texttt{wtmp}) en \texttt{victima2} confirmó que el intruso operó sobre ambas máquinas empleando exactamente el mismo grupo de direcciones IP de acceso conmutado y DSL residencial.

El Listado~\ref{lst:wtmp_victima2_curated} resume los accesos clave registrados en el servidor secundario durante el mes de abril.

\begin{lstlisting}[style=terminal,caption={Extracto representativo de accesos a victima2 desde direcciones IP residenciales.},label={lst:wtmp_victima2_curated}]
duende    pts/4   Sat Apr 26 03:10 - 03:12  (00:02)  dsl-203-0-113-160.isp-acceso.net
duende    pts/3   Tue Apr 22 20:47 - 20:52  (00:04)  dsl-198-51-100-98.isp-acceso.net
duende    pts/0   Thu Apr 17 15:00 - 15:01  (00:01)  dsl-203-0-113-78.isp-acceso.net
duende    pts/2   Wed Apr 16 11:58 - 12:05  (00:07)  dsl-198-51-100-61.isp-acceso.net
zopilote  pts/1   Wed Apr 02 03:05 - 03:06  (00:00)  nodo-acceso.correo-externo.net
\end{lstlisting}

\subsection{Evidencia del Salto Lateral (Pivoting)}

Al contrastar estos registros con los eventos de \texttt{victima1} (analizados en el Capítulo 5), se comprobó la mecánica del salto lateral:
\begin{itemize}
    \item El atacante establecía con frecuencia una sesión SSH inicial desde su conexión externa hacia \texttt{victima2}.
    \item Una vez autenticado en \texttt{victima2}, utilizaba esa máquina como pivote de salto para conectarse localmente hacia \texttt{victima1} a través de la red Fast Ethernet interna.
    \item De este modo, en las bitácoras estándar de \texttt{victima1}, la conexión figuraba originada desde un host local confiable (\texttt{victima2.foo.bar.mx}), intentando con ello desviar la atención de los administradores que únicamente inspeccionaban accesos provenientes de direcciones IP externas.
\end{itemize}

\section{Conclusiones de la Correlación Multisistema}

La auditoría pericial extendida sobre el servidor secundario consolidó las siguientes conclusiones:
\begin{itemize}
    \item La intrusión no constituyó un incidente aislado en \texttt{victima1}, sino una campaña coordinada que abarcó los servidores principales del centro de cómputo.
    \item La presencia de los mismos paquetes de explotación (\texttt{test.tar}, \texttt{ex/}, \texttt{o0o.tar}) y las mismas marcas de hash MD5 en las cuentas \texttt{duende}, \texttt{proa} y \texttt{zena} confirma una autoría común.
    \item El empleo de \texttt{victima2} como nodo de salto (\textit{pivoting}) subraya la importancia forense de no limitar el alcance de una investigación a la máquina donde se detecta la primera alarma, sino auditar de manera integral el perímetro de confianza.
\end{itemize}
